\subsection{Leaf layout and hashing cost}
\label{sec:block-leaves}

With a commitment schedule, the query phase never reads a single block of a committed layer in
isolation. For a committed layer $j_0$ followed by a run of length $m$, every query opens the whole
group of $2^m$ entries of $E_{j_0}$ below the prefix $\omega^{(j_0+m)}$
(Definition of a commitment schedule). The natural Merkle leaf is therefore this group. A tree with
one leaf per pair hashes the same data, but it also builds the $2^{m-1}-1$ internal nodes inside
every group. Their digests are never sent, since the verifier recomputes them from the opened entries.

\paragraph{Hash count.} Let a committed layer have $L$ entries of $\sigma$ bytes, and let
$\kappa(\ell)=\lceil \ell/64\rceil$ be the number of BLAKE3 compressions for an $\ell$-byte input
($\ell\le1024$). With domain tags of $\tau_{\mathrm{leaf}}$, $\tau_{\mathrm{grp}}$ and
$\tau_{\mathrm{node}}$ bytes, the two layouts cost
\[
C_{\mathrm{pair}}=\tfrac{L}{2}\,\kappa(\tau_{\mathrm{leaf}}+2\sigma)+\bigl(\tfrac{L}{2}-1\bigr)\kappa(\tau_{\mathrm{node}}+64),
\qquad
C_{\mathrm{grp}}=\tfrac{L}{2^m}\,\kappa(\tau_{\mathrm{grp}}+2^m\sigma)+\bigl(\tfrac{L}{2^m}-1\bigr)\kappa(\tau_{\mathrm{node}}+64)
\]
compressions. In the implementation $\tau_{\mathrm{leaf}}=\tau_{\mathrm{node}}=22$ and
$\tau_{\mathrm{grp}}=21$, so an internal node costs $\kappa(86)=2$ compressions.

For the first layer ($\sigma=16$) and $m=3$, one group costs $4\cdot1+3\cdot2=10$ compressions with
pair leaves and $\kappa(149)=3$ with a group leaf. Over all committed layers at $n=20$, $k=2$, the
commitment drops from $7{,}489{,}810$ to $3{,}145{,}718$ compressions ($-58\%$) for $m=3$. It drops
by $-35\%$ for $m=2$ and by $-70\%$ for $m=4$.

\paragraph{Proof and verifier.} The opened field elements are the same in both layouts. Each
authentication path is $m-1$ levels shorter, but those levels lie inside the group and were already
recomputed by the verifier. So the proof changes only through the deduplication of shared paths,
which depends on the query positions; in our measurements the change stays within $\pm2\%$. The
verifier hashes one leaf per group instead of $2^{m-1}$ leaves and $2^{m-1}-1$ nodes.

\paragraph{Security.} The proof of Theorem~\ref{thm:schedule} works in the IOP model, where a root is
an oracle to the whole table and the grouping of leaves plays no role. Merkle binding holds for any
fixed public grouping. Group leaves carry their own domain tag, so they cannot be confused with pair
leaves or internal nodes, and the format byte separates the two layouts in the transcript. The
soundness bound, the query count $s$ and the challenge field are unchanged.

\paragraph{Measurements.} Table~\ref{tab:block-leaves} reports the commitment phase alone. It
builds exactly the trees of the prover on the same tables, with the two layouts interleaved in one
process. Table~\ref{tab:block-leaves-e2e} reports full proofs for all four combinations of gate family and
leaf layout.
\begin{table}[h]
\centering
\begin{tabular}{rrrrrr}
\toprule
$n$ & $m$ & pair leaves (ms) & group leaves (ms) & change & compressions\\
\midrule
16 & 3 & 59.2 & 27.9 & $-53\%$ & $468{,}100\to196{,}604$\\
18 & 3 & 236.9 & 111.2 & $-53\%$ & $1{,}872{,}436\to786{,}416$\\
20 & 2 & 1139.9 & 799.5 & $-30\%$ & $9{,}087{,}628\to5{,}941{,}904$\\
20 & 3 & 944.8 & 442.6 & $-53\%$ & $7{,}489{,}810\to3{,}145{,}718$\\
20 & 4 & 864.9 & 302.4 & $-65\%$ & $6{,}850{,}678\to2{,}027{,}234$\\
\bottomrule
\end{tabular}
\caption{Commitment phase, $k=2$, BLAKE3, one thread, Apple M1, median of 20 interleaved
repetitions; the interquartile ranges do not overlap.}
\label{tab:block-leaves}
\end{table}
The time follows the compression count: about $126$\,ns per compression with pair leaves and
$135$--$150$\,ns with group leaves on one thread. The relative changes are the same on an Intel
Xeon at $2.8$\,GHz ($-30\%$, $-56\%$, $-65\%$ at $n=20$).

\begin{table}[h]
\centering
\begin{tabular}{lrrrrrr}
\toprule
& \multicolumn{3}{c}{$m=3$} & \multicolumn{3}{c}{$m=4$}\\
\cmidrule(lr){2-4}\cmidrule(lr){5-7}
gates, leaves & prover & verifier & proof & prover & verifier & proof\\
\midrule
random, pair & 1669 ms & 10.6 ms & 588{,}170 B & 1567 ms & 13.4 ms & 655{,}682 B\\
$(t,t+1)$, pair & 1367 ms & 10.2 ms & 594{,}218 B & 1305 ms & 12.8 ms & 659{,}618 B\\
random, group & 1137 ms & 9.0 ms & 593{,}418 B & 1001 ms & 10.6 ms & 661{,}314 B\\
$(t,t+1)$, group & \textbf{882 ms} & 8.4 ms & 585{,}706 B & \textbf{738 ms} & 9.8 ms & 658{,}658 B\\
\bottomrule
\end{tabular}
\caption{Full proofs at $n=20$, $k=2$, $s=407$, $g=16$, one thread, Apple M1. All four
configurations are built from the same commit and measured in one session, cycling through them in every round (median of 20
runs each). The first row has the cost of v0.9.1 (random gates, pair leaves).}
\label{tab:block-leaves-e2e}
\end{table}

The two optimizations are independent and their gains multiply. Specialized gates
(\S\ref{sec:specialized-gates}) shorten encoding and folding. Group leaves shorten the
commitment. Together they make the single-thread prover $1.89\times$ faster at $m=3$
($1669\to882$\,ms) and $2.12\times$ faster at $m=4$ ($1567\to738$\,ms). The verifier gets faster too,
and the proof size stays within $\pm1\%$ of the baseline for each schedule. With eight threads the
medians are $510\to323$\,ms at $m=3$ and $462\to176$\,ms at $m=4$; the run-to-run spread is larger
there.

The two schedules serve different needs. The \emph{compact} profile $m=3$ gives the smallest proof
($586$\,KB). The \emph{fast} profile $m=4$ gives the fastest prover ($738$\,ms on one thread,
$176$\,ms on eight) for a proof $12\%$ larger. Both have the same soundness bound and query count.

\paragraph{Further saving.} An internal node hashes a $22$-byte tag and $64$ bytes, i.e.\ two
compressions. Moving the tag into the key of a keyed hash leaves a $64$-byte input, i.e.\ one
compression. This would halve the cost of the remaining internal nodes.
